\answer
\begin{enumerate}

%% 5章 問1（5.1節）
\item[5.1節 問1]
$T = 1/f = 1/(100\times10^{3}) = 10\,\mathrm{\mu s}$，
$T_{\mathrm{on}} = DT = 0.3 \times 10 = 3\,\mathrm{\mu s}$。
平均電圧は「$12$\,Vである時間の割合」で決まるから
$D \times 12 = 3.6$\,Vである。
デューティ比の制御がそのまま平均電圧の制御になる。

%% 5章 問2（5.2節）
\item[5.2節 問2]
式\ref{eq5.9}より$D = 12/48 = 0.25$。
オフ期間はダイオードが導通してSの出力側の端子はほぼ接地電位になるから，
Sの両端には入力電圧$V_{\mathrm{in}} = 48$\,Vがかかる。
降圧チョッパのスイッチとダイオードには，
いずれも入力電圧以上の耐圧が必要である。

%% 5章 問3（5.3節）
\item[5.3節 問3]
式\ref{eq5.13}より
$D = 1 - V_{\mathrm{in}}/V_{\mathrm{out}} = 1 - 24/96 = 0.75$。
$D = 0.9$では$V_{\mathrm{out}}/V_{\mathrm{in}} = 1/(1-0.9) = 10$倍が理論上の上限である
（実際には損失のためこの比までは使えない。\ref{sec5.3}節の注意）。

%% 5章 問4（5.4節）
\item[5.4節 問4]
式\ref{eq5.18}より$|V_{\mathrm{out}}|/V_{\mathrm{in}} = D/(1-D)$。
$-5$\,V：$D/(1-D) = 5/12$より$12D = 5 - 5D$，
$D = 5/17 \approx 0.29$。
$-24$\,V：$D/(1-D) = 2$より$D = 2/3 \approx 0.67$。
$D < 0.5$で降圧，$D > 0.5$で昇圧になっていることを確かめてほしい。

%% 5章 問5（5.5節）
\item[5.5節 問5]
式\ref{eq5.21}の右辺の単位は，$D$が無次元であることに注意して
\begin{equation*}
\frac{\mathrm{V}}{\mathrm{Hz}\cdot\mathrm{A}}
= \frac{\mathrm{V}}{(1/\mathrm{s})\cdot\mathrm{A}}
= \frac{\mathrm{V \cdot s}}{\mathrm{A}} = \mathrm{H}
\end{equation*}
となり，左辺（インダクタンス）と一致する。
$1\,\mathrm{H} = 1\,\mathrm{V\cdot s/A}$は式\ref{eq5.2}
（$v_L = L\, di_L/dt$）そのものである。

%% 5章 問6（5.6節）
\item[5.6節 問6]
$V_{\mathrm{out}} = DV_{\mathrm{in}} = 5$\,V。式\ref{eq5.19}より
\begin{equation*}
\mathit{\Delta}I_L = \frac{(10-5)\times0.5}{1\times10^{3}\times30\times10^{-3}}
\approx 83\,\mathrm{mA}
\end{equation*}
式\ref{eq5.24}より
\begin{equation*}
\mathit{\Delta}V_{\mathrm{out}}
= \frac{83\times10^{-3}}{8\times1\times10^{3}\times100\times10^{-6}}
\approx 0.10\,\mathrm{V}
\end{equation*}
出力$5$\,Vに対して約2\,\%のリプルである。

%% 5章 問7（5.7節）
\item[5.7節 問7]
式\ref{eq5.27}より，境界の負荷電流は
$I_{\mathrm{out}} = \mathit{\Delta}I_L/2 = 0.2$\,Aであり，
これ以下でDCMに入る。負荷抵抗では
\begin{equation*}
R = \frac{V_{\mathrm{out}}}{I_{\mathrm{out}}} = \frac{5}{0.2} = 25\,\Omega
\end{equation*}
以上でDCMに入る（軽負荷＝大きい抵抗であることに注意）。

\end{enumerate}
